\documentclass[dvipsnames]{article}
\usepackage{geometry}
\geometry{a3paper, landscape, margin=2in}
\usepackage{url}
\usepackage{multicol}
\usepackage{esint}
\usepackage{amsfonts}
\usepackage{tikz}
\usetikzlibrary{decorations.pathmorphing}
\usepackage{amsmath,amssymb}
\usepackage{enumitem}
\usepackage{colortbl}
\usepackage{mathtools}
\usepackage{ mathrsfs }	
\usepackage{ dsfont }
\usepackage{ gensymb }

\usepackage[activate={true,nocompatibility},final,tracking=true,kerning=true,spacing=true,factor=1100,stretch=10,shrink=10]{microtype}
\makeatletter
\setlist[itemize]{noitemsep, topsep=0pt}
%Template by Drew Ulick https://de.overleaf.com/articles/130-cheat-sheet/ntwtkmpxmgrp
\usepackage[english]{babel}
\usepackage{lmodern}
	
\advance\topmargin-1.2in
\advance\textheight3in
\advance\textwidth3in
\advance\oddsidemargin-1.5in
\advance\evensidemargin-1.5in
\parindent0pt
\parskip2pt
\newcommand{\hr}{\centerline{\rule{3.5in}{1pt}}}
%\colorbox[HTML]{e4e4e4}{\makebox[\textwidth-2\fboxsep][l]{texto}
\begin{document}

\begin{center}{\fontsize{35}{60}{\textcolor{purple}{\textbf{Group Theory Cheat Sheet}}}}\\
\end{center}
\begin{multicols*}{3}

\tikzstyle{mybox} = [draw=purple, fill=white, very thick,
    rectangle, rounded corners, inner sep=10pt, inner ysep=10pt]
\tikzstyle{fancytitle} =[fill=purple, text=white, font=\bfseries]

%------------ Group Definitions ---------------
\begin{tikzpicture}
\node [mybox] (box){%
    \begin{minipage}{0.3\textwidth}
    A group is an ordered pair $(G, *)$ where $G$ is a set and $*$ is a binary operation on $G$ satisfying the following axioms:
    \begin{itemize}
    	\item \textbf{Associativity:} $(a * b) * c = a * (b * c) \forall a, b, c \in G$
    	\item \textbf{Identity:} $\exists e \in G$, called an identity element of $G$, s.t. $\forall a \in G$ we have $a*e = e * a = a$.
    	\item \textbf{Inverse} $\forall a  \in G \exists \ a^{-1} \in G$, called an inverse of $a$, s.t. $a *a^{-1} = a^{-1} * a = e.$
    \end{itemize}
    Closure is guaranteed due to the definition of binary operation.
    
    Identity and inverses are unique.
    \end{minipage}
};
%------------ Group Definitons Header ---------------------
\node[fancytitle, right=10pt] at (box.north west) {Group Axioms};
\end{tikzpicture}
%
%%------------ Some Properties of Groups ---------------
%\begin{tikzpicture}
%\node [mybox] (box){%
%    \begin{minipage}{0.3\textwidth}
%    \begin{itemize}
%    	\item A group G is called \textbf{abelian} if $a * b = b * a \forall a, b \in G$
%    \end{itemize}
%    \textbf{iii. Cancellation property} suppose that a * b = a * c,  $\forall$ a, b, c $\in$ G, $\Rightarrow$ b = c\\
%    \textbf{iv. Uniqueness of Inverse and Identity}\\ • The identity of G is unique\\
%    • $\forall$ a $\in$ G, $a^{-1}$ is uniquely determined\\
%    • $(a^{-1})^{-1}$ = a $\forall$ a $\in$ G\\
%    • $(a * b)^{-1}$ = $(b^{-1})\ *\ (a^{-1})$\\
%    • for any $a_1 , a_2, ... , a_n$ $\in$ G the value of $a_1 * a_2 * · · · *a_n$ is independent of how the expression is bracketed
%    \end{minipage}
%};
%%------------ Mixing Header ---------------------
%\node[fancytitle, right=10pt] at (box.north west) {Some Properties of Groups};
%\end{tikzpicture}

%------------ Some special groups ---------------
\begin{tikzpicture}
\node [mybox] (box){%
    \begin{minipage}{0.3\textwidth}
    	\begin{itemize}
    		\item A group is called \textbf{abelian} if it is commutative.
    		\item The group of all symmetries of a $n-$sided regular polygon is called the \textbf{dihedral group}. It is represented as, \[ D_{2n}= \langle r,s \mid r^{n}=s^{2}=1, rs=sr^{-1} \rangle. \] 
    		\item The group of all bijections on a set on $n$ elements is called the \textbf{symmetric} group denoted as $S_n$.
    		\item The \textbf{Klein-4} group is a group with 4 elements in which each element is a self inverse.
    	\end{itemize}
    \end{minipage}
};
%------------ Some Special Groups Space Header ---------------------
\node[fancytitle, right=10pt] at (box.north west) {Some Special Groups};
\end{tikzpicture}

%------------ Homomorphisms and Isomorphisms ---------------
\begin{tikzpicture}
\node [mybox] (box){%
    \begin{minipage}{0.3\textwidth}
    Let $(G , *)$ and $(H, \circ)$ be groups. A map $\varphi : G \to H,\ s.t.\ \varphi(x * y) = \varphi(x) \circ \varphi(y)\ \forall\ x, y \in G$ is called a \textbf{homomorphism.}\\
    A bijective homomorphism is called an \textbf{isomorphism}.

    \end{minipage}
};
%------------ Homomorphisms and Isomoprhisms Header ---------------------
\node[fancytitle, right=10pt] at (box.north west) {Homomorphisms and Isomorphisms};
\end{tikzpicture}

%------------ Subgroups ---------------------
\begin{tikzpicture}
	\node [mybox] (box){%
		\begin{minipage}{0.3\textwidth}
			For a Group $G$ a subset $H \subseteq G$, is a \textbf{subgroup} of $G$, i.e. $H \leq G$ if it is non empty and a group with the binary operation of $G$ restricted to $H$.
			
			Alternatively, a subset of a group is a subgroup if it is non empty and closed under products and inverses, i.e. for $H \subseteq G$
			\begin{itemize}
				\item $H \neq \emptyset$
				\item $xy^{-1} \in H, \forall x,y \in H$
			\end{itemize}
			
			A subgroup $N$ of $ G$ is called \textbf{normal}, denoted as $N \trianglelefteq G$ if $gng^{-1} \in N, \forall\ g \in G,  n \in N.$
			
		\end{minipage}
	};
	%------------ Subgroups Header ---------------------
	\node[fancytitle, right=10pt] at (box.north west) {Subgroups};
\end{tikzpicture}

%------------ Group Actions ---------------
\begin{tikzpicture}
\node [mybox] (box){%
    \begin{minipage}{0.3\textwidth}
    A \textbf{group action} of a group $(G,*)$ on a set $X$ is a map $\circ:G \times X \to X$ satisfying the following properties,
    \begin{itemize}
    	\item \textbf{Identity:} $e \circ x = x, \forall x \in X$
    	\item \textbf{Compatibility:} $g \circ (h \circ x)=g*h \circ x, \forall g,h \in G, x \in X$
    \end{itemize}
    Alternatively a group action on a set $X$ can be thought of as a homomorphism from $G$ to the symmetric group of $X$.
    
    The \textbf{conjugation action} is a homomorphism $\varphi_x: G \to G$ for some fixed $x \in G$ is defined as $\varphi_x(g)=xgx^{-1}$.

    \end{minipage}
};
%------------ Group Actions Header ---------------------
\node[fancytitle, right=10pt] at (box.north west) {Group Actions};
\end{tikzpicture}


%------------ Stablizer, Kernel of Group Action ---------------------
\begin{tikzpicture}
	\node [mybox] (box){%
		\begin{minipage}{0.3\textwidth}
			If $G$ acts on a set $X$ then we define the following,
			
			\textbf{Orbit} of a $x \in X$ is the set $Gx=\{gx \in X \mid g \in G\}$. Alternatively its the equivalence class induced by the group action.
			
			\textbf{Stablizer} of an element $x \in X$ is the set of elements from the group which leave $x$ fixed, i.e. $G_x=\{g \in G \mid gx=x\}$
			
			\textbf{Kernel} of a group action is the kernel of the associated group homomorphism i.e., $\{g \in G \mid gx=x, \forall x \in X\}$ 
			
			An action is called \textbf{transitive} if it has only one orbit, and is called \textbf{faithful} if its kernel is trivial.
			
			\textbf{Conjugacy classes} of $G$ is the equivalence classes of G when it acts on itself with conjugation. i.e. $\{gag^{-1}\ |\ g \in G \}$
		\end{minipage}
	};
	%------------ Stabl4izer, Kernel of Group Action Header ---------------------
	\node[fancytitle, right=10pt] at (box.north west) {Stablizer, Kernel of Group Action};
\end{tikzpicture}

%------------ Centralizers and Normalizers Content ---------------
\begin{tikzpicture}
\node [mybox] (box){%
    \begin{minipage}{0.3\textwidth}
    \textbf{Centralizer} of $A \subseteq G$ in $G$ is a subset of $G$ defined as $$C_G(A) = \{ g \in G\ |\ gag^{-1} = a\ \forall\ a \in A \}.$$ It is the set of all elements of $G$ which commute with every element of $A$.
    
    \textbf{Center} of $G$ is the subset of G defined as $$Z(G) = \{ g \in G\ |\ gx = xg\ \forall\ x \in G \}.$$ It is the set of elements commutating with all the elements of G. It is the kernel of the conjugation action.
    
    \textbf{Normalizer} of $A$ in $G$ is defined as
    $$N_G(A)=\{ g \in G\ |\ gAg^{-1} = A\}$$ 
    $gAg^{-1}=\{gag^{-1}\ |\ a \in A\}$. Note that $C_G(A) \leq N_G(A).$
   The normalizer of a subset is its stabilizer under conjugation action.
    \end{minipage}
};
%------------ Centralizers and Normalizers Header ---------------------
\node[fancytitle, right=10pt] at (box.north west) {Centralizers and Normalizers};
\end{tikzpicture}


%------------ Cyclic Groups and Cycle Notation Content ---------------
\begin{tikzpicture}
\node [mybox] (box){%
    \begin{minipage}{0.3\textwidth}
    A Group $H$ is \textbf{cyclic} if $\exists\ x \in H\ \text{s.t. } H = \{ x^n\ |\ n \in \mathbb{Z} \}$\\
    For the above case we say $H =\left < x \right >$ and that $H$ is generated by $x$.
    \begin{itemize}
    	\item A cyclic group can have more than one generator.
    	\item All cyclic groups are abelian.
    	\item If $H=\left<x\right>$ then $|H| = |x|,$ if $|H|=n<\infty$ then $x^n=1$
    	\item Any two cyclic groups of the same order are isomorphic.\
    \end{itemize}
   \hrule
   \vspace{0.2cm}
   Two-Line to Cycle notation for permutations\\
   $\left(\begin{matrix}
1 & 2 & 3 & 4 & 5\\ 
2 & 5 & 4 & 3 & 1
\end{matrix}\right)=(125)(34)=(34)(125)=(34)(512)=(15)(25)(34)$\\
Here, the last form is a case of 2-cycle (transposition).
    \end{minipage}
};
%------------ Cyclic Groups Header ---------------------
\node[fancytitle, right=10pt] at (box.north west) {Cyclic Groups and Cycle Notation};
\end{tikzpicture}

%------------ Parity of Permutations and Alternating Groups Content ---------------------
\begin{tikzpicture}
	\node [mybox] (box){%
		\begin{minipage}{0.3\textwidth}
			The parity of any permutation $\sigma$ is given by the parity of the number of its 2-cycles (transpositions).
			\hrule
			\vspace{0.2cm}
			\textbf{Alternating Groups}:\\
			An alternating group is the group of even permutations of a finite set of length $n$. It is denoted by $A_n$ it's order is $\frac{n!}{2}$
			
		\end{minipage}
	};
	%------------ Parity of Permutations and Alternating Groups Header ---------------------
	\node[fancytitle, right=10pt] at (box.north west) {Parity of Permutations and Alternating Groups};
\end{tikzpicture}


%------------ Cosets and Quotient Groups Content ---------------
\begin{tikzpicture}
\node [mybox] (box){%
    \begin{minipage}{0.3\textwidth}
    For any $N\leq G$ and any $g \in G$\\
    • $gN=\{gn\ |\ n\in N\}=\{g, gh_1, gh_2 \dots \}$ and,\\
    • $Ng=\{ng\ |\ n\in N\}=\{g, h_1g, h_2g \dots \}$
    are called a left coset and a right coset respectively.\\
    \hrule
    \vspace{0.2cm}
    For a Group $G\ \text{and}\ N \trianglelefteq G$, the \textbf{quotient group} of N in G (i.e. G/N), is the set of cosets of N in G.
    
    The definition of a normal subgroup is the same as left and right cosets being equal.
    
    \end{minipage}
};
%------------ Cosets and Quotient Groups Header ---------------------
\node[fancytitle, right=10pt] at (box.north west) {Cosets and Quotient Groups};
\end{tikzpicture}

%------------ Lagrange's Theorem and some results ---------------
\begin{tikzpicture}
\node [mybox] (box){%
    \begin{minipage}{0.3\textwidth}
    \textbf{Lagrange's Theorem}: For a finite group $G$ and $H \leq G$,\\
    • The order of $H$ divides the order of $G$, and,\\
    • The number of left cosets of $H$ in $G$ equals $\frac{|G|}{|H|}$\\
    \hrule
    \vspace{0.2cm}
    \textbf{Some important results}\\
    • If $G$ is a finite group and $x \in G$, then the order of $x$ divides the order of $G$, and $x^{|G|}=e\ \forall\ x \in G$\\
    • If $G$ is a group of prime order, then $G$ is cyclic
    \end{minipage}
};
%------------ Lagrange's Theorem and some results ---------------------
\node[fancytitle, right=10pt] at (box.north west) {Lagrange's Theorem and some results};
\end{tikzpicture}
\
%------------ Cauchy's Theorem Content ---------------
\begin{tikzpicture}
\node [mybox] (box){%
    \begin{minipage}{0.3\textwidth}
    \textbf{Cauchy's Theorem}: If $G$ is a finite group and $p$ is a prime dividing $|G|$ then $G$ has an element of order $p$.
    \end{minipage}
};
%------------ Cauchy's Theorems Header ---------------------
\node[fancytitle, right=10pt] at (box.north west) {Cauchy's Theorem};
\end{tikzpicture}
%------------ The Isomorphism Theorems Content ---------------------
\begin{tikzpicture}
\node [mybox] (box){%
    \begin{minipage}{0.3\textwidth}
	\textbf{First Isomorphism Theorem:}\\
	If $\varphi:G \to H$ is a homomorphism of groups. Then $\mathrm{ker} \varphi \trianglelefteq G$ and,
	$G/\mathrm{ker}\varphi \cong \varphi(G)$.
	
	\textbf{Second Isomorphism Theorem:}\\
	For a group $G$ with, $A, B \leq G$ and, $A \trianglelefteq N_G(B)$.
Then $AB \leq G, \\ B\trianglelefteq AB, A \cap B \trianglelefteq A$ and, $AB/B \cong A/ A \cap B$\\
	\textbf{The Third Isomorphism Theorem:}\\
	For a group $G$ with, $H, K \trianglelefteq G$ and, $H\leq K$. Then $K/H \trianglelefteq G/H$ and, $\frac{G/H}{K/H}\cong G/K$
	\end{minipage}
};
%------------ The Isomorphism Theorems Header ---------------------
\node[fancytitle, right=10pt] at (box.north west) {The Isomorphism Theorems};
\end{tikzpicture}


%------------ Orbit–stabilizer Theorem ---------------
\begin{tikzpicture}
	\node [mybox] (box){%
		\begin{minipage}{0.3\textwidth}
			\textbf{Orbit–stabilizer Theorem:}\\
			Let a finite group $G$ act on a set $X$. For $x \in X$ denote its orbit by $Gx$
			and its stabilizer by $G_x$. Then the map
			$G/G_x \to Gx$, $gG_x \mapsto gx$, is a bijection, so
			$|Gx| = [G : G_x]$ and hence $|Gx|\,|G_x| = |G|$.
		\end{minipage}
	};
	\node[fancytitle, right=10pt] at (box.north west) {Orbit–stabilizer Theorem};
\end{tikzpicture}

%------------ Orbit decomposition and class equation ---------------
\begin{tikzpicture}
	\node [mybox] (box){%
		\begin{minipage}{0.3\textwidth}
			For a finite group $G$ acting on a finite set $X$ we have the
			\textbf{orbit decomposition}
			$X = \bigsqcup_i Gx_i$, where $x_i$ run over orbit representatives, so
			$|X| = \sum_i |Gx_i| = \sum_i [G : G_{x_i}]$.
			For the conjugation action of $G$ on itself, the orbits are conjugacy classes
			and the fixed points are $Z(G)$. The resulting formula
			$$|G| = |Z(G)| + \sum_i [G : C_G(x_i)]$$
			(with $x_i$ running over non-central conjugacy classes) is the
			\textbf{class equation} of $G$.
		\end{minipage}
	};
	\node[fancytitle, right=10pt] at (box.north west) {Orbit decomposition and class equation};
\end{tikzpicture}
%------------ p-group actions and fixed points ---------------
\begin{tikzpicture}
	\node [mybox] (box){%
		\begin{minipage}{0.3\textwidth}
			Let $G$ be a finite $p$-group acting on a finite set $X$. Write
			$\mathrm{Fix}(X) = \{x \in X \mid gx = x \ \forall g \in G\}$. Then
			$$|X| \equiv |\mathrm{Fix}(X)| \pmod p.$$
			In particular, if $p \nmid |X|$, the action has a fixed point.
			Applying this to the conjugation action of a $p$-group on itself gives
			the class equation
			$|G| = |Z(G)| + p k$, so $Z(G) \ne 1$ and more generally every finite
			$p$-group has nontrivial center.
		\end{minipage}
	};
	\node[fancytitle, right=10pt] at (box.north west) {p-group actions and fixed points};
\end{tikzpicture}


%------------ Cayley's Theorem Content ---------------
\begin{tikzpicture}
\node [mybox] (box){%
    \begin{minipage}{0.3\textwidth}
    \textbf{Cayley's Theorem:}\\
    Every group is isomorphic to a subgroup of some symmetric group. If $G$ is a group of order $n$, then G is isomoprhic to a subgroup of $S_n$
    \end{minipage}
};
%------------ Cayley's Theorem Header ---------------------
'\node[fancytitle, right=10pt] at (box.north west) {Cayley's Theorem};
\end{tikzpicture}

%------------ Automorphisms ---------------
\begin{tikzpicture}
	\node [mybox] (box){%
		\begin{minipage}{0.3\textwidth}
			An \textbf{automorphism} of $G$ is an isomorphism $G \to G$. The set
			$\mathrm{Aut}(G)$ is a group under composition.
			For $g \in G$, the map $x \mapsto gxg^{-1}$ is an \textbf{inner automorphism}.
			The inner automorphisms form a normal subgroup
			$\mathrm{Inn}(G) \le \mathrm{Aut}(G)$ with
			$\mathrm{Inn}(G) \cong G/Z(G)$. The quotient
			$\mathrm{Out}(G) = \mathrm{Aut}(G)/\mathrm{Inn}(G)$ is the group of
			\textbf{outer automorphisms}.
			Some facts:\\
			$\bullet$ $\mathrm{Aut}(\mathbb{Z}/n\mathbb{Z}) \cong (\mathbb{Z}/n\mathbb{Z})^\times$.\\
			$\bullet$ For $n \ne 2,6$, every automorphism of $S_n$ is inner, so
			$\mathrm{Aut}(S_n) = \mathrm{Inn}(S_n)$.\\
			$\bullet$ $S_6$ has a nontrivial outer automorphism:
			$\mathrm{Out}(S_6) \ne 1$.
		\end{minipage}
	};
	\node[fancytitle, right=10pt] at (box.north west) {Automorphisms};
\end{tikzpicture}

%------------ p-groups and Sylow p-groups ---------------
\begin{tikzpicture}
	\node [mybox] (box){%
		\begin{minipage}{0.3\textwidth}
			A \textbf{p-group} is a finite group of order $p^a$ for some prime $p$ and $a \ge 1$.
			Subgroups of $G$ which are $p$-groups are called $p$-subgroups.
			If $|G| = p^a m$ with $p \nmid m$, a \textbf{Sylow p-subgroup} of $G$ is a
			subgroup of order $p^a$. The set of Sylow $p$-subgroups is denoted $Syl_p(G)$.
		\end{minipage}
	};
	\node[fancytitle, right=10pt] at (box.north west) {p-groups and Sylow p-groups};
\end{tikzpicture}
%------------ Sylow Theorems ---------------
\begin{tikzpicture}
	\node [mybox] (box){%
		\begin{minipage}{0.3\textwidth}
			Let $|G| = p^a m$ with $p \nmid m$.
			
			\textbf{First Sylow Theorem:}
			$G$ has a subgroup of order $p^a$ (a Sylow $p$-subgroup).
			
			\textbf{Second Sylow Theorem:} Any two Sylow $p$-subgroups of $G$ are conjugate.

			\textbf{Third Sylow Theorem:}
			If $n_p$ is the number of Sylow $p$-subgroups, then
			$n_p \equiv 1 \pmod p$ and $n_p \mid m$.

			\textbf{(“Fourth” Sylow consequence):}
			A Sylow $p$-subgroup $P$ is normal in $G$ iff $n_p = 1$ (equivalently,
			$P$ is the unique Sylow $p$-subgroup).
		\end{minipage}
	};
	\node[fancytitle, right=10pt] at (box.north west) {Sylow Theorems};
\end{tikzpicture}

%------------ Sylow consequences ---------------
\begin{tikzpicture}
	\node [mybox] (box){%
		\begin{minipage}{0.3\textwidth}
			Let $G$ be a finite group and $p$ a prime.
			
			$\bullet$ Every $p$-subgroup of $G$ is contained in some Sylow $p$-subgroup.\\
			$\bullet$ If $H \le G$ and $P$ is a Sylow $p$-subgroup of $G$, then
			$P \cap H$ is a Sylow $p$-subgroup of $H$.\\
			$\bullet$ If $H \le G$ and $[G:H]$ is coprime to $p$, then any
			Sylow $p$-subgroup of $H$ is a Sylow $p$-subgroup of $G$.\\
			$\bullet$ If $[G:H]$ is a power of $p$, then $H$ contains a Sylow
			$p$-subgroup of $G$.
		\end{minipage}
	};
	\node[fancytitle, right=10pt] at (box.north west) {Sylow consequences};
\end{tikzpicture}
%------------ Frattini subgroup and Frattini argument ---------------
\begin{tikzpicture}
	\node [mybox] (box){%
		\begin{minipage}{0.3\textwidth}
			The \textbf{Frattini subgroup} $\Phi(G)$ of a group $G$ is the
			intersection of all maximal subgroups of $G$ (for a finite $p$-group,
			the intersection of its maximal subgroups). It is characteristic in $G$.
			
			For a finite $p$-group $G$, the quotient $G/\Phi(G)$ is an elementary
			abelian $p$-group; its dimension over $\mathbb{F}_p$ equals the minimal
			number of generators of $G$.
			
			\textbf{Frattini argument:} If $N \trianglelefteq G$ and $P$ is a Sylow
			$p$-subgroup of $N$, then $G = N_G(P)\,N$.
		\end{minipage}
	};
	\node[fancytitle, right=10pt] at (box.north west) {Frattini subgroup and Frattini argument};
\end{tikzpicture}

%------------ Commutators Content ---------------
\begin{tikzpicture}
	\node [mybox] (box){%
		\begin{minipage}{0.3\textwidth}
			For a group $G$ and $x,y \in G$ call $[x,y]=x^{-1}y^{-1}xy$ the \textbf{commutator} of $x$ and $y$.
			
			For subsets $A,B \subseteq G$ define the group generated by its commutators as $[A,B]=\langle [a,b] \mid a\in A, b \in B\rangle$.
			
			Similarly $G'=[G,G]$ is the subgroup of $G$ generated by all commutators, its called the commutator subgroup of $G$, or its \textbf{derived subgroup}.	
			
			The following are some useful properties of commutators,
			\begin{itemize}
				\item $xy=yx \iff [x,y]=1$
				\item $H \trianglelefteq G \iff [H,G] \leq H$
				\item $G/G'$ is abelian and its the largest abelian quotient of $G.$ It is called the \textbf{abelianization} of $G$.
				\item Any homomorphism from $G$ to an abelian group $A$ factors through $G/G'$ (via the natural map $G \to G/G'$).
			
			\end{itemize}
		\end{minipage}
	};
	%------------ Commutators Header ---------------------
	'\node[fancytitle, right=10pt] at (box.north west) {Commutators};
	
\end{tikzpicture}

%------------ Direct products Content ---------------
\begin{tikzpicture}
	\node [mybox] (box){%
		\begin{minipage}{0.3\textwidth}
			The \textbf{direct product} $G_1 \times G_2 \times \cdots $of groups $G_1, G_2, \dots$ is set of sequences $(g_1,g_2, \cdots)$ with $g_i \in G_i$ and operation $*$ defined component wise.
			
			If $H, K$ are normal subgroups of $G$ with $H \cap K =1$ then $HK \cong H \times K$
		\end{minipage}
	};
	%------------ Direct products Header ---------------------
	'\node[fancytitle, right=10pt] at (box.north west) {Direct products};
	
\end{tikzpicture}
%------------ Semidirect products ---------------
\begin{tikzpicture}
	\node [mybox] (box){%
		\begin{minipage}{0.3\textwidth}
			\textbf{External semidirect product:}\\
			Given groups $H, K$ and a homomorphism $\varphi : K \to \mathrm{Aut}(H)$,
			the semidirect product $H \rtimes_\varphi K$ is the set $H \times K$ with
			multiplication
			$$(h_1,k_1)(h_2,k_2) = (h_1\,\varphi(k_1)(h_2),\,k_1k_2).$$[2pt]
			\textbf{Internal semidirect product:}\\
			A group $G$ is an internal semidirect product of subgroups $H,K$ if\\
			$\bullet$ $H \trianglelefteq G$, \quad $\bullet$ $HK = G$, \quad $\bullet$ $H \cap K = 1$.\\
			Then the map $H \rtimes_\varphi K \to G$, $(h,k) \mapsto hk$, is an isomorphism
			for a suitable $\varphi : K \to \mathrm{Aut}(H)$ given by conjugation.
			Different $\varphi$ that differ by inner automorphisms of $H$ give isomorphic
			semidirect products; only the outer class of $\varphi$ matters.
		\end{minipage}
	};
	\node[fancytitle, right=10pt] at (box.north west) {Semidirect products};
\end{tikzpicture}


%------------ Exact sequences and extensions ---------------
\begin{tikzpicture}
	\node [mybox] (box){%
		\begin{minipage}{0.3\textwidth}
			A \textbf{short exact sequence} of groups is
			$$1 \to N \xrightarrow{i} G \xrightarrow{\pi} Q \to 1,$$
			where $i$ is injective, $\pi$ is surjective and $\mathrm{im}(i) = \ker(\pi)$.\\
			This encodes that $N \trianglelefteq G$ and $G/N \cong Q$; one says that
			$G$ is an \textbf{extension} of $Q$ by $N$.
			The sequence \textbf{splits} if there is a homomorphism
			$s : Q \to G$ with $\pi \circ s = \mathrm{id}_Q$. In this case
			$G \cong N \rtimes Q$ is a semidirect product.
		\end{minipage}
	};
	\node[fancytitle, right=10pt] at (box.north west) {Exact sequences and extensions};
\end{tikzpicture}


%------------ Fundamental theorem of finitely generated abelian groups ---------------
\begin{tikzpicture}
	\node [mybox] (box){%
		\begin{minipage}{0.3\textwidth}
			Every finitely generated abelian group $A$ is isomorphic to a direct product
			of cyclic groups. Equivalently, there is a decomposition (noncanonically)
			$$A \cong \mathbb{Z}^r \times \mathbb{Z}/n_1\mathbb{Z} \times \cdots \times \mathbb{Z}/n_k\mathbb{Z},$$
			with $r \ge 0$ and $n_1 \mid n_2 \mid \cdots \mid n_k$. The integers $r$ and
			$n_i$ are uniquely determined.
		\end{minipage}
	};
	\node[fancytitle, right=10pt] at (box.north west) {Finitely generated abelian groups};
\end{tikzpicture}



%------------ Schur-Zassenhaus Theorem Content ---------------
\begin{tikzpicture}
	\node [mybox] (box){%
		\begin{minipage}{0.3\textwidth}
			\textbf{Schur-Zassenhaus Theorem:}\\
			Let $G$ be a finite group and $N \trianglelefteq G$ with $\gcd(|N|,|G/N|) = 1$.
			Then there exists a subgroup $H \leq G$ such that $|H| = |G/N|$ and
			$G = N \rtimes H$ with $N \cap H = 1$. Any two such complements $H$ are conjugate in $G$.
		\end{minipage}
	};
	%------------ Schur-Zassenhaus Theorem Header ---------------------
	\node[fancytitle, right=10pt] at (box.north west) {Schur-Zassenhaus Theorem};
\end{tikzpicture}


%------------ Simple groups and Composition Series Content ---------------
\begin{tikzpicture}
	\node [mybox] (box){%
		\begin{minipage}{0.3\textwidth}
			A group $G$ is called \textbf{simple} if $|G|>1$ and its only normal subgroups are $1$ and $G$.
			
			A sequence of subgroups as follows (called a subnormal sequence),\[ 1=N_0 \trianglelefteq N_1 \trianglelefteq \cdots \trianglelefteq N_{k-1} \leq N_k = G \] is called a \textbf{composition series} if $N_{i+1}/N_i$ is simple, for $i=0,\dots k-1$. The quotient groups are called composition factors.
		\end{minipage}
	};
	%------------ Simple groups and Composition Series Header ---------------------
	'\node[fancytitle, right=10pt] at (box.north west) {Simple groups and Composition Series};
\end{tikzpicture}

%------------ Zassenhaus lemma and Schreier refinement ---------------
\begin{tikzpicture}
	\node [mybox] (box){%
		\begin{minipage}{0.3\textwidth}
			\textbf{Zassenhaus (butterfly) lemma:}\\
			Let $G$ be a group with subgroups $A,C$ and normal subgroups
			$B \trianglelefteq A$, $D \trianglelefteq C$. Then there is an isomorphism
			$$\frac{(A \cap C)B}{(A \cap D)B} \cong
			\frac{(A \cap C)D}{(B \cap C)D}.$$
			This is visualized by the “butterfly” diagram in the subgroup lattice.
			
			\textbf{Schreier refinement theorem:}
			Any two subnormal series of a group $G$ admit equivalent refinements
			(their factor groups can be refined so that the multisets of factors are
			isomorphic up to order). Jordan–Hölder is the special case for
			composition series.
		\end{minipage}
	};
	\node[fancytitle, right=10pt] at (box.north west) {Zassenhaus lemma and Schreier refinement};
\end{tikzpicture}


%------------ Jordan–Hölder theorem ---------------
\begin{tikzpicture}
	\node [mybox] (box){%
		\begin{minipage}{0.3\textwidth}
			A \textbf{composition series} of a nontrivial finite group $G$ is a subnormal
			series
			$$1 = N_0 \trianglelefteq N_1 \trianglelefteq \cdots \trianglelefteq N_k = G$$
			with each factor $N_{i+1}/N_i$ simple. These $N_{i+1}/N_i$ are the
			\textbf{composition factors}.
			
			\textbf{Jordan–Hölder theorem:}
			Every finite nontrivial group has a composition series, and any two
			composition series have the same multiset of composition factors,
			up to isomorphism and order.
		\end{minipage}
	};
	\node[fancytitle, right=10pt] at (box.north west) {Jordan–Hölder theorem};
\end{tikzpicture}



%------------ Examples of simple groups ---------------
\begin{tikzpicture}
	\node [mybox] (box){%
		\begin{minipage}{0.3\textwidth}
			Basic examples of finite simple groups:
			\begin{itemize}
				\item Simple abelian groups are exactly the cyclic groups $C_p$ of prime order $p$.
				\item For $n \ge 5$, the alternating group $A_n$ is non-abelian and simple.
				\item Projective special linear groups: for a field $F$,
				$SL_n(F) = \{A \in GL_n(F) \mid \det A = 1\}$, 
				$PSL_n(F) = SL_n(F)/Z(SL_n(F))$.
				Over a finite field $\mathbb{F}_q$, $PSL_n(\mathbb{F}_q)$ is simple for $n \ge 2$
				except $(n,q) = (2,2)$ and $(2,3)$; in particular $PSL_2(q)$ is simple for all $q \ge 4$.
			\end{itemize}
		\end{minipage}
	};
	%------------ Examples of simple groups Header ---------------------
	\node[fancytitle, right=10pt] at (box.north west) {Examples of simple groups};
\end{tikzpicture}
%------------ Free groups and presentations ---------------
\begin{tikzpicture}
	\node [mybox] (box){%
		\begin{minipage}{0.3\textwidth}
			A \textbf{free group} on a set $X$ is a group $F(X)$ with an
			injection $X \to F(X)$ such that any map $X \to G$ into a group $G$
			extends uniquely to a homomorphism $F(X) \to G$ (universal property).
			Elements are reduced words in the generators and their inverses.
			
			A \textbf{presentation} of a group is
			$G = \langle X \mid R \rangle$, the quotient of $F(X)$ by the normal
			subgroup generated by a set of relators $R$. Every group has some
			presentation; if both $X$ and $R$ are finite, $G$ is
			\textbf{finitely presented}.
			
			(Nielsen–Schreier) Every subgroup of a free group is free.
		\end{minipage}
	};
	\node[fancytitle, right=10pt] at (box.north west) {Free groups and presentations};
\end{tikzpicture}


%------------ Solvable groups Content ---------------
\begin{tikzpicture}
	\node [mybox] (box){%
		\begin{minipage}{0.3\textwidth}
			A group $G$ is called \textbf{solvable} if it admits a subnormal series
			$$1 = G_0 \trianglelefteq G_1 \trianglelefteq \cdots \trianglelefteq G_n = G$$
			with each factor $G_{i+1}/G_i$ abelian.\\[2pt]
			Equivalently, the \textbf{derived series}
			$$G^{(0)} = G,\quad G^{(1)} = [G,G],\quad G^{(i+1)} = [G^{(i)},G^{(i)}]$$
			terminates: $G^{(n)} = 1$ for some $n$.\\[4pt]
			Basic properties:\\
			$\bullet$ If $H \le G$ and $G$ is solvable, then $H$ is solvable.\\
			$\bullet$ If $\varphi : G \to K$ is a homomorphism and $G$ is solvable,
			then $\varphi(G)$ and $G/\ker\varphi$ are solvable.\\
			$\bullet$ If $N \trianglelefteq G$ and both $N$ and $G/N$ are solvable,
			then $G$ is solvable (extensions of solvable groups are solvable).\\
			$\bullet$ A finite direct product of solvable groups is solvable.\\[4pt]
			For finite groups, $G$ is solvable $\iff$ all composition factors of $G$
			are cyclic of prime order. Burnside: if $|G| = p^a q^b$ ($p,q$ primes),
			then $G$ is solvable. Feit–Thompson: every finite group of odd order
			is solvable.
		\end{minipage}
	};
	%------------ Solvable groups Header ---------------------
	\node[fancytitle, right=10pt] at (box.north west) {Solvable groups};
\end{tikzpicture}

%------------ Nilpotent groups Content ---------------
\begin{tikzpicture}
	\node [mybox] (box){%
		\begin{minipage}{0.3\textwidth}
			A group $G$ is called \textbf{nilpotent} if its \textbf{lower central series}
			$$G_0 = G,\quad G_{i+1} = [G_i,G]$$
			terminates: $G_n = 1$ for some $n \ge 0$. The least such $n$ is the
			nilpotency class.\\[4pt]
			Equivalent descriptions:\\
			$\bullet$ $G$ has a finite \textbf{central series}
			$$1 = G_0 \trianglelefteq G_1 \trianglelefteq \cdots \trianglelefteq G_n = G$$
			with each factor $G_{i+1}/G_i \subseteq Z(G/G_i).$\\
			$\bullet$ (Upper central series) There is a chain
			$$1 = Z_0 \subseteq Z_1 \subseteq \cdots \subseteq Z_c = G$$
			with $Z_{i+1}/Z_i = Z(G/Z_i)$ for all $i$; this terminates iff $G$ is nilpotent.\\
			$\bullet$ For finite $G$: $G$ is nilpotent $\iff$ $G$ is the (internal)
			direct product of its Sylow $p$-subgroups $\iff$ every Sylow subgroup is normal.\\
			$\bullet$ (Normalizer condition, for finite $G$) $G$ is nilpotent
			$\iff$ every proper subgroup $H < G$ satisfies $H < N_G(H)$.\\[4pt]
			Closure properties:\\
			$\bullet$ Subgroups of nilpotent groups are nilpotent.\\
			$\bullet$ Homomorphic images and quotients of nilpotent groups are nilpotent.\\
			$\bullet$ A finite direct product of nilpotent groups is nilpotent.\\[4pt]
			In particular, every finite $p$-group is nilpotent, and every nilpotent
			group is solvable.
		\end{minipage}
	};
	%------------ Nilpotent groups Header ---------------------
	\node[fancytitle, right=10pt] at (box.north west) {Nilpotent groups};
\end{tikzpicture}

%------------ Linear representations ---------------
\begin{tikzpicture}
	\node[mybox](box){%
		\begin{minipage}{0.3\textwidth}
			Let $G$ be a finite group over $\mathbb{C}$.
			\begin{itemize}
				\item A (linear) representation of $G$ is a homomorphism
				$\rho : G \to GL(V)$, where $V$ is a finite-dimensional
				complex vector space.
				\item The degree of $\rho$ is $\dim_{\mathbb{C}} V$.
				\item Choosing a basis of $V$ identifies $\rho$ with
				$\rho : G \to GL_n(\mathbb{C})$ (matrix representation).
				\item Two representations $(V,\rho)$ and $(W,\sigma)$ are
				equivalent if there exists a linear isomorphism
				$T : V \to W$ with $T\rho(g) = \sigma(g)T$ for all $g$.
				\item Finite-dimensional representations of $G$ over $\mathbb{C}$
				are the same as finite-dimensional left modules over
				the group algebra $\mathbb{C}G$.
			\end{itemize}
		\end{minipage}
	};
	\node[fancytitle,right=10pt] at (box.north west){Linear representations};
\end{tikzpicture}

%------------ Subrepresentations and irreducibility ---------------
\begin{tikzpicture}
	\node[mybox](box){%
		\begin{minipage}{0.3\textwidth}
			Let $(V,\rho)$ be a representation of $G$.
			\begin{itemize}
				\item A subspace $W \subseteq V$ is $G$-invariant if
				$\rho(g)W \subseteq W$ for all $g \in G$.
				Then $W$ is a subrepresentation.
				\item The quotient $V/W$ becomes a representation via
				$\rho(g)(v+W) = \rho(g)v + W$.
				\item A nonzero representation is irreducible if it has no proper
				nonzero subrepresentation.
				\item A representation is reducible if it is not irreducible.
				\item A representation is completely reducible if it is a direct
				sum of irreducible subrepresentations.
				\item Direct sum: if $(V_i,\rho_i)$ are representations, then
				$\bigoplus_i V_i$ with $g\cdot(v_1,\dots,v_k)
				= (\rho_1(g)v_1,\dots,\rho_k(g)v_k)$ is a representation.
			\end{itemize}
		\end{minipage}
	};
	\node[fancytitle,right=10pt] at (box.north west){Subrepresentations and irreducibility};
\end{tikzpicture}

%------------ Basic constructions and examples ---------------
\begin{tikzpicture}
	\node[mybox](box){%
		\begin{minipage}{0.3\textwidth}
			\begin{itemize}
				\item Trivial representation: $V = \mathbb{C}$, $\rho(g) = 1$.
				\item Sign representation of $S_n$: $V = \mathbb{C}$,
				$\rho(\sigma) = \mathrm{sgn}(\sigma)$.
				\item Permutation representation: if $G$ acts on a finite set $X$,
				let $V = \mathbb{C}^X$ with basis $\{e_x\}$ and
				$g\cdot e_x = e_{gx}$.
				\item Regular representation: $V = \mathbb{C}G$ with basis
				$\{e_g\}$ and $h\cdot e_g = e_{hg}$.
				\item Tensor product: for $V,W$, define $g\cdot(v\otimes w)
				= \rho_V(g)v \otimes \rho_W(g)w$.
				\item Dual representation: $V^* = \mathrm{Hom}_{\mathbb{C}}(V,\mathbb{C})$
				with $(g\cdot f)(v) = f(\rho(g^{-1})v)$.
			\end{itemize}
		\end{minipage}
	};
	\node[fancytitle,right=10pt] at (box.north west){Basic constructions and examples};
\end{tikzpicture}

%------------ Intertwiners and Schur's Lemma ---------------
\begin{tikzpicture}
	\node[mybox](box){%
		\begin{minipage}{0.3\textwidth}
			Let $(V,\rho)$ and $(W,\sigma)$ be representations.
			\begin{itemize}
				\item A homomorphism of representations (intertwiner) is a
				linear map $T : V \to W$ with $T\rho(g) = \sigma(g)T$ for all $g$.
				\item Denote by $\mathrm{Hom}_G(V,W)$ the space of intertwiners.
				Put $\mathrm{End}_G(V) = \mathrm{Hom}_G(V,V)$.
				\item Schur's Lemma (1): if $V,W$ are irreducible, any
				$T \in \mathrm{Hom}_G(V,W)$ is either $0$ or an isomorphism.
				\item Schur's Lemma (2): if $V$ is irreducible over $\mathbb{C}$,
				then $\mathrm{End}_G(V)$ consists of scalar multiples of the
				identity.
			\end{itemize}
		\end{minipage}
	};
	\node[fancytitle,right=10pt] at (box.north west){Intertwiners and Schur's Lemma};
\end{tikzpicture}

%------------ Unitary representations and Maschke's theorem ---------------
\begin{tikzpicture}
	\node[mybox](box){%
		\begin{minipage}{0.3\textwidth}
			\begin{itemize}
				\item A Hermitian inner product $(\cdot,\cdot)$ on $V$ is
				$G$-invariant if $(\rho(g)v,\rho(g)w) = (v,w)$ for all $g$.
				Then $\rho(g)$ is unitary.
				\item For any inner product $(\cdot,\cdot)$ on $V$, define
				$(v,w)_G = \frac{1}{|G|}\sum_{g\in G}(\rho(g)v,\rho(g)w)$.
				Then $(\cdot,\cdot)_G$ is $G$-invariant.
				\item Over $\mathbb{C}$, every finite group representation is
				equivalent to a unitary representation.
				\item Maschke's Theorem: if $\mathrm{char}(\mathbb{C})=0$
				(or more generally $\mathrm{char}(F)\nmid |G|$), every
				representation of $G$ is completely reducible.
				\item Equivalently, for every $G$-invariant subspace
				$W\subseteq V$ there exists a $G$-invariant complement
				$W'$ with $V = W \oplus W'$.
			\end{itemize}
		\end{minipage}
	};
	\node[fancytitle,right=10pt] at (box.north west){Unitary representations and Maschke's theorem};
\end{tikzpicture}

%------------ Characters: definition and basic facts ---------------
\begin{tikzpicture}
	\node[mybox](box){%
		\begin{minipage}{0.3\textwidth}
			Let $(V,\rho)$ be a representation.
			\begin{itemize}
				\item The character of $V$ is $\chi_V : G \to \mathbb{C}$ given by
				$\chi_V(g) = \mathrm{tr}(\rho(g))$.
				\item $\chi_V(1) = \dim_{\mathbb{C}} V$.
				\item If $V\cong W$ as representations, then $\chi_V = \chi_W$.
				\item For irreducible representations over $\mathbb{C}$,
				equality of characters implies equivalence.
				\item $\chi_V$ is constant on conjugacy classes
				(class function).
				\item For direct sums and tensor products:
				$\chi_{V\oplus W} = \chi_V + \chi_W$, and
				$\chi_{V\otimes W} = \chi_V\chi_W$ (pointwise).
				\item If $V$ is unitary, then $\chi_{V^*}(g)
				= \overline{\chi_V(g)}$.
				\item For a permutation representation from an action on $X$,
				$\chi(g)$ equals the number of fixed points of $g$ in $X$.
			\end{itemize}
		\end{minipage}
	};
	\node[fancytitle,right=10pt] at (box.north west){Characters: definition and basic facts};
\end{tikzpicture}

%------------ Inner product and orthogonality of characters ---------------
\begin{tikzpicture}
	\node[mybox](box){%
		\begin{minipage}{0.3\textwidth}
			Let $\mathrm{CF}(G)$ be the space of complex class functions on $G$.
			\begin{itemize}
				\item Inner product on $\mathrm{CF}(G)$:
				$\langle \varphi,\psi\rangle
				= \frac{1}{|G|}\sum_{g\in G}\varphi(g)\overline{\psi(g)}$.
				\item If $\chi_V,\chi_W$ are characters, then
				$\langle \chi_V,\chi_W\rangle$ is a nonnegative integer.
				It equals the multiplicity of $W$ in $V$.
				\item Orthogonality: if $\chi_i,\chi_j$ are irreducible
				characters, then $\langle \chi_i,\chi_j\rangle = \delta_{ij}$.
				\item A character $\chi$ is irreducible iff
				$\langle \chi,\chi\rangle = 1$.
				\item If $V \cong \bigoplus_i m_i V_i$ with $V_i$ irreducible,
				then $\chi_V = \sum_i m_i\chi_i$ and
				$m_i = \langle \chi_V,\chi_i\rangle$.
			\end{itemize}
		\end{minipage}
	};
	\node[fancytitle,right=10pt] at (box.north west){Inner product and orthogonality of characters};
\end{tikzpicture}

%------------ Main character formulas ---------------
\begin{tikzpicture}
	\node[mybox](box){%
		\begin{minipage}{0.3\textwidth}
			Let $C_1,\dots,C_h$ be the conjugacy classes of $G$ and
			$\chi_1,\dots,\chi_h$ the irreducible characters.
			\begin{itemize}
				\item The number of irreducible characters equals the number
				of conjugacy classes.
				\item The irreducible characters form an orthonormal basis of
				$\mathrm{CF}(G)$ with respect to the inner product above.
				\item Let $d_i = \chi_i(1)$ be the degrees of irreducible
				representations. Then
				$\sum_i d_i^2 = |G|$.
				\item Row orthogonality: for $i\neq j$,
				$\sum_{g\in G}\chi_i(g)\overline{\chi_j(g)} = 0$, and
				for $i=j$, the sum is $|G|$.
				\item Column orthogonality: if $g,h$ are representatives of
				different conjugacy classes, then
				$\sum_i \chi_i(g)\overline{\chi_i(h)} = 0$; if $g,h$ lie in
				the same class $C$, then
				$\sum_i \chi_i(g)\overline{\chi_i(h)} = |G|/|C|$.
			\end{itemize}
		\end{minipage}
	};
	\node[fancytitle,right=10pt] at (box.north west){Main character formulas};
\end{tikzpicture}

%------------ Regular representation and group algebra ---------------
\begin{tikzpicture}
	\node[mybox](box){%
		\begin{minipage}{0.3\textwidth}
			Let $R$ be the regular representation on $V = \mathbb{C}G$.
			\begin{itemize}
				\item The character of $R$ satisfies
				$\chi_R(1) = |G|$ and $\chi_R(g) = 0$ for $g\neq 1$.
				\item If $V_i$ are the irreducible representations with
				degrees $d_i$, then
				$R \cong \bigoplus_i d_i V_i$.
				\item Comparing characters gives
				$\sum_i d_i^2 = |G|$ and shows that $V_i$ occurs in $R$
				with multiplicity $d_i$.
				\item As an algebra,
				$\mathbb{C}G \cong \prod_i M_{d_i}(\mathbb{C})$.
				\item There exist central idempotents $e_i \in \mathbb{C}G$
				with $1 = \sum_i e_i$, $e_ie_j = 0$ for $i\neq j$, and
				$e_i\mathbb{C}G \cong M_{d_i}(\mathbb{C})$.
			\end{itemize}
		\end{minipage}
	};
	\node[fancytitle,right=10pt] at (box.north west){Regular representation and group algebra};
\end{tikzpicture}

%------------ Permutation characters and Burnside's lemma ---------------
\begin{tikzpicture}
	\node[mybox](box){%
		\begin{minipage}{0.3\textwidth}
			Let $G$ act on a finite set $X$ and consider the permutation
			representation on $V = \mathbb{C}^X$.
			\begin{itemize}
				\item The character satisfies
				$\chi(g) = |\{x\in X \mid gx = x\}|$.
				\item The number of $G$-orbits in $X$ is
				$\frac{1}{|G|}\sum_{g\in G}\chi(g)$.
				This is the Cauchy–Frobenius (Burnside) lemma.
				\item For the conjugation action of $G$ on itself,
				permutation characters encode the sizes of conjugacy classes
				and centralizers.
			\end{itemize}
		\end{minipage}
	};
	\node[fancytitle,right=10pt] at (box.north west){Permutation characters and Burnside's lemma};
\end{tikzpicture}

%------------ Restriction, induction and Frobenius reciprocity ---------------
\begin{tikzpicture}
	\node[mybox](box){%
		\begin{minipage}{0.3\textwidth}
			Let $H\leq G$ be a subgroup.
			\begin{itemize}
				\item Restriction: a representation $V$ of $G$ becomes a
				representation of $H$ by restricting $\rho$ to $H$.
				For characters: $(\mathrm{Res}^G_H\chi)(h)=\chi(h)$ for $h\in H$.
				\item Induction: a representation $W$ of $H$ has an induced
				representation $\mathrm{Ind}^G_H W$ of $G$ with
				$\dim(\mathrm{Ind}^G_H W) = [G:H]\dim W$.
				\item Frobenius reciprocity (characters):
				if $\theta$ is a character of $H$ and $\chi$ a character of $G$, then
				$\langle \mathrm{Ind}^G_H\theta,\chi\rangle_G
				= \langle \theta,\mathrm{Res}^G_H\chi\rangle_H$.
			\end{itemize}
		\end{minipage}
	};
	\node[fancytitle,right=10pt] at (box.north west){Restriction, induction and Frobenius reciprocity};
\end{tikzpicture}

%------------ Irreducible representations of abelian groups ---------------
\begin{tikzpicture}
	\node[mybox](box){%
		\begin{minipage}{0.3\textwidth}
			Let $G$ be a finite abelian group.
			\begin{itemize}
				\item Every irreducible representation of $G$ over $\mathbb{C}$
				is $1$-dimensional.
				\item Irreducible representations are group homomorphisms
				$\chi : G \to \mathbb{C}^\times$.
				\item There are exactly $|G|$ distinct irreducible characters;
				they form the dual group $\widehat{G}$ under pointwise product.
				\item The regular representation of $G$ decomposes as the direct
				sum of all irreducible characters, each with multiplicity $1$.
				\item The character table of $G$ is a square matrix of size
				$|G|\times |G|$ and is a discrete Fourier transform on $G$.
			\end{itemize}
		\end{minipage}
	};
	\node[fancytitle,right=10pt] at (box.north west){Irreducible representations of abelian groups};
\end{tikzpicture}

%------------ Example: representations of S_3 ---------------
\begin{tikzpicture}
	\node[mybox](box){%
		\begin{minipage}{0.3\textwidth}
			For $S_3$:
			\begin{itemize}
				\item Conjugacy classes:
				identity $(1)$, transpositions $(12)$, $3$-cycles $(123)$.
				Class sizes: $1,3,2$.
				\item There are $3$ irreducible representations:
				\begin{itemize}
					\item Trivial: degree $1$, character values $(1,1,1)$.
					\item Sign: degree $1$, character values $(1,-1,1)$.
					\item Standard $2$-dimensional: degree $2$, character
					values $(2,0,-1)$.
				\end{itemize}
				\item Degrees satisfy
				$1^2 + 1^2 + 2^2 = 6 = |S_3|$.
				\item The regular representation of $S_3$ decomposes as
				$R \cong 1 \oplus 1_{\mathrm{sgn}} \oplus 2\cdot V_{\mathrm{std}}$.
			\end{itemize}
		\end{minipage}
	};
	\node[fancytitle,right=10pt] at (box.north west){Example: representations of $S_3$};
\end{tikzpicture}



\end{multicols*}
\end{document}


© 2016 GitHub, Inc. Terms Privacy Security Status Help